\documentclass[11pt,a4paper]{article}

\usepackage[utf8]{inputenc}
\usepackage[T1]{fontenc}
\usepackage[portuguese]{babel}
\usepackage{lmodern}
\usepackage{microtype}
\usepackage[a4paper,margin=2.7cm]{geometry}
\usepackage{booktabs}
\usepackage{longtable}
\usepackage{enumitem}
\usepackage[hidelinks]{hyperref}
\usepackage[round,authoryear]{natbib}

\hypersetup{pdftitle={Evolução temporal do debate parlamentar português},pdfauthor={A preencher}}
\title{Evolução temporal do debate parlamentar português\\\large Documento de trabalho metodológico}
\author{[Nome do autor]\\[0.25em][Mestrado / Programa]\\[Instituição]}
\date{\today}

\begin{document}
\maketitle

\begin{abstract}
Este documento regista, de forma cumulativa e auditável, as decisões metodológicas de um estudo sobre a evolução temática do debate na Assembleia da República portuguesa. A investigação usa o corpus ParlaMint-PT e prepara intervenções parlamentares para uma análise temporal de tópicos baseada em embeddings e no enquadramento BERTilda. Esta versão consolida o problema de investigação, a revisão de literatura essencial, a proveniência e o pré-processamento dos dados e a análise exploratória que antecede a modelação. As decisões ainda não validadas empiricamente são identificadas explicitamente como provisórias.
\end{abstract}

\noindent\textbf{Palavras-chave:} debate parlamentar; modelação de tópicos; evolução temporal; embeddings; ParlaMint; BERTilda.

\tableofcontents
\newpage

\section{Introdução}

Os registos parlamentares permitem estudar, em grande escala, a atenção política, a competição entre partidos e a organização do debate ao longo do tempo. Porém, a frequência de palavras não é uma medida suficiente da estrutura temática: termos diferentes podem remeter para o mesmo problema político e um mesmo termo pode mudar de enquadramento. O objetivo desta investigação é caracterizar e seguir a evolução de temas no debate parlamentar português, preservando a unidade empírica da intervenção e usando a dimensão temporal para comparar tópicos entre períodos.

O corpus de partida é o ParlaMint-PT, parte da família ParlaMint de corpora parlamentares comparáveis e codificados em TEI. Estes recursos combinam transcrições com metadados sobre sessões, pessoas oradoras, partidos e outras propriedades institucionais \citep{erjavec2023parlamint}. A disponibilidade simultânea de texto e metadados permite que a análise seja interpretada à luz de mudanças de governo, composição partidária e ritmo parlamentar.

O problema metodológico central não é apenas descobrir tópicos em cada período. É determinar quando dois tópicos observados em períodos adjacentes representam continuidade, quando um tópico se divide ou converge com outro e quando a aparente mudança decorre de baixa atividade ou de composição desequilibrada do corpus. O enquadramento BERTilda trata precisamente esta tarefa como alinhamento entre tópicos de janelas adjacentes e classificação explícita de eventos de ciclo de vida \citep{oliveirafigueira2026bertilda}.

\subsection{Questões de investigação}

\begin{enumerate}[label=\textbf{Q\arabic*.}]
  \item Como evolui a atenção temática no debate parlamentar português entre 2015 e março de 2024?
  \item Que continuidades, aparecimentos, desaparecimentos, divisões e convergências temáticas podem ser identificados entre períodos consecutivos?
  \item Em que medida a interpretação dessa evolução depende da composição partidária, do estatuto das pessoas oradoras e da cadência parlamentar?
  \item Que granularidade temporal oferece comparações interpretáveis sem confundir recessos e alterações do ritmo parlamentar com mudança temática?
\end{enumerate}

As questões Q1--Q3 definem o objetivo substantivo. A Q4 é instrumental: uma janela temporal não altera a construção dos embeddings por intervenção, mas condiciona a descoberta e o alinhamento de tópicos feitos depois.

\section{Revisão de literatura}

\subsection{Corpora parlamentares e análise computacional do discurso}

Os debates parlamentares são fontes particularmente úteis para análise computacional porque combinam continuidade temporal, registo institucional e metadados sociopolíticos. O projeto ParlaMint uniformiza recursos parlamentares de vários países, organiza-os em TEI e disponibiliza formatos derivados para análise, reduzindo barreiras de reutilização e comparação \citep{erjavec2023parlamint}. Esta infraestrutura não elimina problemas de cobertura, de transcrição ou de classificação, mas permite tornar explícita a proveniência das variáveis usadas.

O estudo de \citet{roswevers2025epistemic} ilustra o potencial de uma abordagem temporal e estrutural. Os autores modelam debates da câmara baixa neerlandesa, constroem redes de associação entre tópicos em períodos semestrais e estudam a especialização discursiva através de densidade, modularidade e comunidades. O seu desenho não é um modelo a copiar diretamente: usa LDA, PMI e Louvain, enquanto este projeto prevê um método baseado em embeddings. Ainda assim, mostra que a granularidade deve ser justificada pelo ritmo institucional e que resultados temporais exigem interpretação substantiva.

\subsection{De modelos probabilísticos a tópicos baseados em embeddings}

O modelo LDA representa documentos como misturas de tópicos e tópicos como distribuições sobre palavras \citep{blei2003lda}. As extensões dinâmicas, como o Dynamic Topic Model, introduzem dependência temporal nas distribuições temáticas \citep{bleilafferty2006dtm}; modelos mais recentes, como o Dynamic Embedded Topic Model, evoluem representações em espaços de embeddings \citep{dieng2019detm}. Estas abordagens são úteis quando se pretende impor uma trajetória temporal suave, mas essa imposição pode ocultar mudanças abruptas ou tornar menos auditável a identificação de divisões e convergências.

Em contraste, a modelação de tópicos baseada em embeddings pode descobrir tópicos independentemente em cada período e só depois estabelecer correspondências temporais. BERTopic combina embeddings de documentos, agrupamento e uma variante class-based TF--IDF para obter representações interpretáveis dos tópicos \citep{grootendorst2022bertopic}. BERTilda usa esta lógica de descoberta por instantâneos, constrói um grafo temporal entre tópicos de janelas adjacentes e combina similaridade semântica com cobertura bidirecional de documentos para rotular continuidades, divisões, convergências, desaparecimentos e relações ambíguas \citep{oliveirafigueira2026bertilda}.

\subsection{Embeddings de intervenções}

Modelos Transformer contextualizam a representação de uma palavra pelo seu contexto \citep{devlin2019bert}. Para comparação semântica e agrupamento, são frequentes representações de frase ou documento que possam ser comparadas no espaço vetorial; Sentence-BERT foi proposto precisamente para produzir embeddings comparáveis eficientemente por similaridade \citep{reimersgurevych2019sbert}. O presente projeto iniciou um piloto com \textit{neuralmind/bert-base-portuguese-cased} (BERTimbau), um modelo pré-treinado para português brasileiro \citep{souza2020bertimbau}, e mean pooling normalizado.

Esta última escolha é uma \emph{configuração experimental}, não uma conclusão: BERTimbau com mean pooling não é automaticamente um codificador Sentence-BERT afinado para similaridade. A adequação ao português europeu, à linguagem parlamentar e à tarefa de agrupamento terá de ser avaliada contra alternativas e por inspeção humana das relações temáticas. Esta ressalva evita transformar conveniência computacional em justificação científica.

\subsection{Avaliação e interpretabilidade}

A qualidade de um modelo de tópicos não pode ser deduzida apenas de uma métrica automática. Medidas de coerência são úteis, mas podem divergir de julgamentos humanos \citep{chang2009reading,newman2010automatic}. Para uma análise temporal, acrescentam-se perguntas próprias: se as relações entre períodos são estáveis, se os eventos de ciclo de vida são plausíveis e se a interpretação resiste à inspeção de documentos. BERTilda propõe combinar diagnósticos quantitativos de qualidade e estabilidade com validação humana de relações temporais \citep{oliveirafigueira2026bertilda}. Esta investigação adotará o mesmo princípio: métricas para triagem e comparação, e amostras rastreáveis de intervenções para validar interpretações.

\section{Dados e pré-processamento}

\subsection{Fonte, unidade de análise e âmbito temporal}

A unidade de análise é a intervenção parlamentar, não uma data nem uma janela temporal. Cada linha liga um identificador de intervenção ao seu texto, data, sessão e metadados da pessoa oradora e da afiliação política. O corpus final cobre 7 de janeiro de 2015 a 13 de março de 2024. Assim, 2024 é um ano incompleto e não deve ser comparado diretamente com anos completos.

Foram usados os ficheiros tabulares derivados do ParlaMint-PT para juntar texto e metadados. A estrutura TEI foi também usada para consultar afiliações institucionais de pessoas oradoras, em particular mandatos parlamentares e funções no Governo. Esta combinação preserva a tabela de análise enquanto permite verificar classificações temporais.

\subsection{Filtragem institucional e integração}

\begin{enumerate}
  \item Foram carregados os ficheiros \texttt{*-meta.tsv} e convertida a variável \texttt{Date}; o ano foi derivado da data.
  \item A coluna \texttt{Agenda} foi retirada por não contribuir para o objetivo atual. Foram também retiradas \texttt{Body} e \texttt{Lang}, esta última sem variação útil no corpus.
  \item Foram excluídas intervenções com \texttt{Speaker\_role = Convidado}. Posteriormente, foram excluídas intervenções da presidência e da Mesa (\texttt{Presidente de mesa} e \texttt{Mesa}). Esta decisão reduz texto procedimental, mas não garante a eliminação de todo o discurso procedimental.
  \item Os textos foram lidos dos ficheiros correspondentes \texttt{.txt} e unidos aos metadados pelo identificador \texttt{ID}, com verificação de unicidade e de ausência de textos em falta.
  \item \texttt{Party\_status} foi renomeada para \texttt{bloco\_politico}; \texttt{Coalition} foi mapeada para Governo e \texttt{Opposition} para Oposição. O rótulo descreve posição formal relativamente ao executivo, não apoio parlamentar.
\end{enumerate}

O registo histórico do notebook indica 248\,573 intervenções antes da filtragem de presidência/Mesa e 179\,711 após essa filtragem. Este último é o ponto de partida do processamento textual.

\subsection{Duas representações textuais}

Foram criadas duas versões alinhadas pelos mesmos identificadores e metadados. A separação evita aplicar ao modelo contextual uma normalização pensada apenas para descritores lexicais.

\begin{longtable}{p{0.22\linewidth}p{0.34\linewidth}p{0.36\linewidth}}
\toprule
Representação & Transformações & Uso previsto \\
\midrule
\endhead
\texttt{texto\_embeddings} & Remove apenas marcas editoriais delimitadas por \texttt{[[...]]} e normaliza espaços. Preserva maiúsculas, acentos, pontuação e números. & Embeddings contextuais e experiências de agrupamento. \\
\addlinespace
\texttt{texto\_lexical} & Minúsculas, tokenização e lematização com \texttt{pt\_core\_news\_sm}; remove pontuação, stopwords e termos procedimentais provisórios. Preserva números e as negações \emph{não}, \emph{nem}, \emph{nunca}, \emph{jamais} e \emph{sem}. & Descritores, inspeção lexical e comparação exploratória com a versão contextual. \\
\bottomrule
\end{longtable}

O processamento lexical usa spaCy \citep{honnibal2020spacy}. A lista de termos procedimentais inclui formas de tratamento, referências à presidência e marcas de interação, por exemplo \emph{senhor}, \emph{deputada}, \emph{aplausos} e \emph{protestos}. A lista é deliberadamente provisória: a frequência de uma palavra não demonstra que ela seja procedural, e remover termos sem inspeção de contexto pode eliminar conteúdo político substantivo.

Antes da lematização, foram excluídas intervenções com menos de dez palavras no texto minimamente limpo. Das 179\,711 intervenções de entrada, 60\,298 (33,6\%) foram removidas por este limiar, ficando 119\,413. Dezasseis intervenções sem tokens lexicais após a normalização foram retiradas de ambas as versões para manter os identificadores alinhados. Cada corpus final contém, por isso, 119\,397 intervenções.

\subsection{Registo do piloto de embeddings}

Foi executado um piloto nos anos 2015--2017, com 36\,258 intervenções por versão textual. A configuração usou BERTimbau cased, comprimento máximo de 512 tokens, mean pooling, normalização L2, lotes de quatro documentos em MPS e ficheiros de saída em blocos de 512 intervenções. Cada execução inclui uma assinatura SHA-256 dos identificadores, textos e parâmetros, para impedir a mistura silenciosa de corpus ou configurações.

Este piloto demonstra viabilidade técnica e reprodutibilidade de armazenamento; não constitui ainda uma avaliação da qualidade semântica dos embeddings nem uma seleção definitiva de representação textual ou modelo.

\section{Análise exploratória de dados}

\subsection{Objetivo e princípios}

A análise exploratória antecede a modelação para identificar assimetrias de volume, composição e atividade que possam enviesar a leitura de tópicos. Não serve para escolher uma janela temporal antes de construir embeddings: os embeddings são produzidos ao nível da intervenção. A dimensão temporal é um diagnóstico adicional para a descoberta e comparação de tópicos posterior.

\subsection{Perfil do corpus final}

\begin{table}[htbp]
\centering
\caption{Perfil do corpus lexical usado no EDA.}
\label{tab:perfil-corpus}
\begin{tabular}{lr}
\toprule
Indicador & Valor \\
\midrule
Intervenções & 119\,397 \\
Período & 2015-01-07 a 2024-03-13 \\
Intervenções classificadas como Oposição & 82\,707 \\
Intervenções classificadas como Governo & 35\,012 \\
Intervenções sem rótulo Governo/Oposição & 1\,678 \\
\bottomrule
\end{tabular}
\end{table}

O volume anual não é uniforme. Em particular, 2023 contém 23\,731 intervenções, enquanto 2024 contém apenas 1\,478 por terminar em março. As diferenças de volume não devem ser interpretadas automaticamente como diferenças de saliência temática: podem refletir calendário parlamentar, duração das sessões, procedimentos ou cobertura.

\subsection{Pessoas oradoras e estatuto institucional}

O EDA cruza \texttt{Speaker\_MP} com a lista TEI de pessoas e períodos de afiliação ao Governo. Isto permite distinguir, para cada intervenção rotulada como Governo ou Oposição, deputadas/deputados, membros do Governo e outros não deputados. A distinção é temporal: uma mesma pessoa pode ser deputada num momento e exercer função governativa noutro. Por isso, não se infere o estatuto numa intervenção apenas a partir de uma biografia agregada.

A classificação Governo/Oposição é especialmente útil para contextualizar resultados, mas não deve ser confundida com identidade partidária. Um partido pode mudar de posição formal entre governos; análises por partido e por bloco político respondem, portanto, a perguntas diferentes.

\subsection{Composição partidária e concentração}

No corpus final, PS e PSD concentram o maior número de intervenções, seguidos por PCP, CDS-PP e BE. CH, IL e Livre têm presença concentrada nos anos mais recentes, de acordo com a sua entrada ou representação parlamentar. Há ainda partidos e independentes com baixo volume, para os quais médias de extensão ou tópicos específicos podem ser instáveis.

Esta assimetria tem uma consequência direta: tópicos globais tendem a refletir desproporcionalmente o texto dos grupos com maior produção. A análise posterior deverá reportar cobertura por partido e usar limiares explícitos quando interpretar resultados desagregados. Não se deve concluir que a ausência de um tópico num partido com pouca atividade equivale a ausência de posição política.

\subsection{Atividade temporal, composição e sazonalidade}

Foram descritas séries semanais, quinzenais e mensais de intervenções, palavras e extensão textual. A série semanal tem variabilidade elevada; períodos sem intervenções refletem recessos ou interrupções parlamentares, e não documentos sem texto. A agregação mensal suaviza essa variação, mas esta observação não escolhe, por si só, a janela de modelação.

Para tornar a comparabilidade mensurável, foram calculados por período observado: intervenções, palavras, dias ativos, número e número efetivo de pessoas oradoras e partidos, HHI, quota dos cinco oradores/partidos mais ativos e proporção Governo/Oposição. O número efetivo de categorias é definido por
\[
N_{\mathrm{efetivo}} = \frac{1}{\sum_i p_i^2},
\]
e o HHI por $\sum_i p_i^2$. Estas medidas descrevem concentração e diversidade; não medem qualidade temática. A sua função é sinalizar períodos onde uma comparação posterior pode estar dominada por poucas pessoas ou partidos.

Por fim, a sazonalidade agrega todos os anos por mês do calendário e mostra intervenções e dias ativos. Esta análise identifica o ritmo parlamentar, mas os meses posteriores a março ficam subcontados por causa da cobertura parcial de 2024. Uma versão comparável da sazonalidade deverá, numa revisão futura, normalizar por número de anos observados por mês.

\section{Método previsto e decisões em aberto}

\subsection{Pipeline analítico previsto}

\begin{enumerate}
  \item gerar embeddings por intervenção para uma ou mais representações textuais;
  \item descobrir tópicos independentemente em cada janela temporal, mantendo uma configuração comparável;
  \item representar tópicos e ligar tópicos de janelas adjacentes por similaridade semântica e cobertura bidirecional;
  \item construir um grafo temporal e classificar relações como continuidade, aparecimento, desaparecimento, divisão, convergência ou ambiguidade;
  \item avaliar tópicos e relações através de métricas, amostras de intervenções e validação humana.
\end{enumerate}

As janelas serão aplicadas apenas a partir do passo 2. A escolha entre semana, quinzena, mês, semestre ou outra granularidade deve ser sustentada pela atividade observada, diversidade de participantes, interpretabilidade dos tópicos e estabilidade das relações, não apenas pelo número médio de documentos.

\subsection{Decisões a validar antes de resultados substantivos}

\begin{longtable}{p{0.31\linewidth}p{0.27\linewidth}p{0.32\linewidth}}
\toprule
Decisão & Estado atual & Evidência necessária \\
\midrule
\endhead
Representação textual contextual versus lexical & Em comparação & Qualidade e diversidade dos tópicos; inspeção de amostras; impacto nas relações temporais. \\
\addlinespace
Modelo de embeddings & Piloto com BERTimbau + mean pooling & Comparação com modelo afinado para similaridade ou alternativa adequada ao português europeu; validação humana de vizinhos semânticos. \\
\addlinespace
Granularidade temporal & Sem decisão final & Atividade, concentração, composição, coerência e estabilidade de relações por janela. \\
\addlinespace
Parâmetros de agrupamento / BERTopic & Não fixados & Sensibilidade a hiperparâmetros, proporção de ruído, coerência e auditabilidade. \\
\addlinespace
Limiar e regras de relações temporais & Não fixados & Precisão em amostra anotada e análise de casos limítrofes. \\
\bottomrule
\end{longtable}

\section{Limitações conhecidas}

O estudo cobre texto parlamentar, não a totalidade da atividade política. A filtragem de presidência e Mesa reduz ruído procedimental, mas não remove necessariamente todos os atos procedimentais. A lematização e a lista de termos removidos podem introduzir erros, sobretudo por o modelo linguístico e BERTimbau terem sido desenvolvidos predominantemente para português brasileiro. Finalmente, a representação de Governo/Oposição é institucional e temporal, mas não mede ideologia, disciplina de voto ou apoio parlamentar informal.

Estas limitações não invalidam a análise; delimitam as inferências aceitáveis e definem testes de robustez para as próximas versões.

\section{Próximos passos}

\begin{enumerate}
  \item Rever manualmente a lista de termos procedimentais e uma amostra de lemas em português europeu.
  \item Comparar embeddings contextuais e lexicais com uma tarefa de vizinhança semântica e inspeção humana.
  \item Definir uma grelha de experiências temporalmente comparável, variando uma dimensão de cada vez.
  \item Estabelecer critérios de avaliação de tópicos e de relações temporais antes de interpretar resultados políticos.
  \item Exportar figuras e tabelas do EDA para inclusão nesta secção, com legendas que indiquem cobertura temporal e filtros aplicados.
\end{enumerate}

\bibliographystyle{plainnat}
\bibliography{references}
\end{document}
